\section{Approach: Log-to-Code Retrieval Plus Incident Memory}
\label{sec:approach}

Our approach, \sys{}, follows one observation: software failures recur.
When a new issue arrives, some past issue in the same project has often
failed in a similar way, and its fix commit names the code that was at
fault. \sys{} therefore combines two kinds of evidence. The first is
\emph{log-to-code retrieval}: matching the report's text and logs against
the snapshot's files. The second is \emph{incident memory}: retrieving
similar resolved incidents and letting their fix commits vote for files.
This section describes both, the firewall that keeps the memory honest,
and the fusion that produces the final ranking.

\subsection{Overview}
\label{sec:approach:overview}

The input matches the challenge: an issue report with its logs, and the
code snapshot. The output is a ranked list of the snapshot's source
files, the affected module derived from that list, and, for every ranked
file, the evidence that put it there. Four signals each produce a ranked
list over the snapshot's non-test source files. A weighted fusion merges
the four lists into one. All components run on CPU, use no training, and
are deterministic: the same input always yields the same ranking.

\subsection{Log-to-code signals}
\label{sec:approach:signals}

\textbf{Trace mentions (S1).} The most direct evidence is a file named in
the report itself. S1 finds explicit mentions: file names with an
extension (\texttt{KafkaChannel.java}, \texttt{row0mysql.cc:123}) and
class names inside stack-trace frames (\texttt{at
org.apache.kafka...}). Mentioned files are ranked by where they first
appear in the text, since reporters tend to paste the decisive trace
early. If the text names no file explicitly, S1 falls back to counting
occurrences of bare file-name stems. S1 is deliberately simple because it
doubles as a baseline: it is the automated form of an engineer grepping
the report.

\textbf{Path retrieval (S2).} Path names carry meaning: a report about a
stale SSL factory shares words with
\texttt{security/ssl/SslFactory.java}. S2 scores every file path against
the report with BM25, a standard lexical ranking function, after
splitting identifiers on case and underscores.

\textbf{Content retrieval (S3).} The strongest textual evidence sits
inside the files: identifiers, string literals, and the logging
statements that produced the reported log lines. Scoring every file's
content in every snapshot would be slow and, worse, would tempt a global
index that leaks information across time. S3 therefore scores a
per-instance candidate pool: the files found by S1, the top 200 files
from S2, and the files proposed by memory (below). Term statistics are
computed inside this pool only, so nothing from any other snapshot or any
later version influences the score.

\subsection{The incident memory signal}
\label{sec:approach:memory}

\textbf{Memory retrieval and voting (S4).} S4 searches the project's
memory pool (\cref{sec:benchmark:memory}) for past incidents whose
summary and description read like the current report, again with BM25
over an inverted index. The twenty most similar admissible incidents then
vote: an incident at similarity rank $r$ adds $1/r$ to every gold file of
its own fix commit, provided that file still exists in the current
snapshot. The intuition is the on-call engineer's: ``this looks like the
broker bug from last spring, and that one ended in
\texttt{KafkaChannel}.'' Unlike S1 to S3, S4 can point at a file that no
log line and no trace ever names, which is exactly what the hard stratum
requires.

\subsection{The memory firewall}
\label{sec:approach:firewall}

A memory signal evaluated on historical data can cheat in ways that are
invisible if they are not controlled. We exclude a past incident from
voting on a query when any of the following holds, and we count every
exclusion:

\begin{enumerate}
\item it is the query issue itself;
\item it belongs to the query's duplicate cluster, so it describes the
  same incident under another key;
\item it shares a fix commit with the query, which happens when one
  commit closes several issues;
\item it was resolved after the query was created. The memory an
  engineer could actually search at report time contains only incidents
  resolved before the report existed.
\end{enumerate}

Rule 4 is strict and it costs real signal: in our runs it blocks roughly
two thirds of candidate votes. We accept the cost because the remaining
signal is one a deployed system would truly have had.

\subsection{Fusion, modules, and evidence}
\label{sec:approach:fusion}

\textbf{Weighted fusion.} The four lists are merged by reciprocal rank
fusion: file $f$ receives $\sum_i w_i / (60 + \mathrm{rank}_i(f))$, where
$\mathrm{rank}_i(f)$ is $f$'s position in signal $i$'s list. The weights
$w_i$ are the only tuned parameters in \sys{}.

\textbf{Conditioning on an observable regime.} The right weights depend
on the report. When the text names files explicitly, S1 is nearly always
right and should dominate. When it names none, S1 is noise and memory
matters more. The easy/hard strata of \cref{sec:benchmark:strata} cannot
be used for this decision, because they are defined through the gold
answer. We instead condition on a property that is visible at inference
time: whether S1 found at least one explicit file mention. Weights are
tuned separately for the two regimes, by grid search over
$\{0, 0.5, 1, 2\}$ per signal, on a development split only
(\cref{sec:setup:protocol}). The ablation without memory is tuned the
same way, so the comparison in \cref{sec:results} is fair to both sides.

\textbf{Module and evidence outputs.} The affected module is read off the
file ranking: files collapse to their leading path components (one to
three, fixed per project since two of the four keep all code under
\texttt{src/}), and a module's rank is the rank of its first file. For
the explanation output, every ranked file carries its provenance: which
signals ranked it, the matching log lines or mentions for S1 to S3, and
for S4 the past incidents that voted for it with their issue keys. An
engineer can follow these links directly to the prior fix.
